% !PRINTED 699
\weberpage{699}
If now $(O, B') > m$, then in $B'$ one could find a polygon
$\mathfrak{M}\mathfrak{B}$ divisible by $\mathfrak{M}$ (\S\,192, 5.), and then
\[
\mathfrak{A}' \mathfrak{M} \mathfrak{B} = \mathfrak{A} \mathfrak{B} = \mathfrak{W}
\]
would be a complete polygon of the first kind; $\mathfrak{A}$ itself would
therefore be of the first kind, contrary to the assumption. Hence
$(O, B') = (A', W) = m$ and
\weq{15}{(O, A) = a - p + 1.}

So here too the \gesperrt{Riemann-Roch theorem} holds in the
form 3.

If $A$ consists of an isolated polygon $\mathfrak{M}$, then $(O, A) = 1$,
hence $a = p$, and from this it follows:

6. \gesperrt{An isolated polygon of the second kind must\
be a $p$-gon, and conversely every $p$-gon
of the second kind is an isolated polygon.}

7. \gesperrt{If a point $\mathfrak{P}$ divides the divisor of an improper
 class $A$ of the second kind, and if $A = \mathfrak{P} A'$,
then $A'$ must be of the first kind.}

For by the theorem (6) and (15) we have:
\[
\begin{aligned}
(O, A) &= a - p + 1, \\
(O, A') &= a - p + (A', W), \\
(O, A) &= (O, A'),
\end{aligned}
\]
hence $(A', W) = 1$.

\begin{center}
\textbf{\S\,200. Differentials of the second and third kind.}
\end{center}
We can now represent invariantly the condition for a differential:
\weq{1}{dJ = \frac{\mathfrak{A}}{\mathfrak{B}},}
which we expressed in \S\,197, (6) with reference to a variable $z$ in the
form:
\[
\mathfrak{U}^2 \mathfrak{A} \sim \mathfrak{Z} \mathfrak{B},
\]
by using the complete polygons of the first kind, according to
which
\[
\mathfrak{U}^2 \mathfrak{W} \sim \mathfrak{Z}
\]
holds:
\[
\mathfrak{A} \sim \mathfrak{W} \mathfrak{B},
\]
or, by introducing the classes:
%CONT>>
